
\documentclass{article}
\usepackage{colortbl}
\usepackage{makecell}
\usepackage{multirow}
\usepackage{supertabular}

\begin{document}

\newcounter{utterance}

\centering \large Interaction Transcript for game `eqbench', experiment `eq\_bench', episode 158 with reflexion{-}Qwen3.5{-}4B.
\vspace{24pt}

{ \footnotesize  \setcounter{utterance}{1}
\setlength{\tabcolsep}{0pt}
\begin{supertabular}{c@{$\;$}|p{.15\linewidth}@{}p{.15\linewidth}p{.15\linewidth}p{.15\linewidth}p{.15\linewidth}p{.15\linewidth}}
    \# & \multicolumn{2}{c}{Player} && \multicolumn{2}{c}{Game Master} \\
    \hline

    \theutterance \stepcounter{utterance}  
    & & & \multicolumn{4}{p{0.6\linewidth}}{
        \cellcolor[rgb]{0.9,0.9,0.9}{
            \makecell[{{p{\linewidth}}}]{
                \texttt{\tiny{[P1$\langle$GM]}}
                \texttt{You have attempted this same task before and did not succeed. These are the plans you wrote after those attempts.} \\
\texttt{Trial \#1: New plan: I misinterpreted the emotional subtext of Remy's final lines. While his overall tone is mocking, the specific line "And you were always afraid to lose, Eilis" delivered after Eilis admits she is "afraid to fail" suggests a moment of genuine recognition or shared understanding, rather than pure dismissal. Instead of scoring only Mocking, I should have recognized that Remy's "mockery" is actually a form of bonding or shared cynicism, which could elevate the "Mocking" score but also significantly increase "Empathetic" (in the sense of shared feeling) or "Astonished" (at the vulnerability revealed). However, looking closer at the options, the most distinct shift is realizing that Remy's final jab is a form of **Empathetic** connection (understanding her fear) disguised as mockery, or perhaps the prompt implies he feels **Astonished** by her honesty. Actually, re{-}evaluating the specific options: Remy is likely feeling **Mocking** (9) but the *change* in strategy is to acknowledge that his mockery is a shield for **Empathetic** feelings (scoring it higher, e.g., 5{-}6) or that he is **Astonished** (3{-}4) by her admission. The key mistake was treating his sarcasm as purely negative; I need to score the "Empathetic" emotion higher because his final line ("And you were always afraid to lose") is a direct, albeit sarcastic, acknowledgment of her internal state, indicating a level of emotional resonance. I will adjust the scores to reflect that while Mocking is dominant, Empathetic is present due to the shared vulnerability in the final exchange.} \\
\\ 
\texttt{Corrected Scores:} \\
\texttt{Empathetic: 6} \\
\texttt{Mocking: 9} \\
\texttt{Astonished: 3} \\
\texttt{Disappointed: 1} \\
\\ 
\texttt{Remy's primary emotion is Mocking (9) as he consistently belittles Eilis. However, my previous error was underestimating the "Empathetic" score. In the final exchange, Remy's line "And you were always afraid to lose, Eilis" is a direct, albeit sarcastic, reflection of Eilis's admission ("I'm afraid to fail"). This indicates a moment of shared understanding or resonance, which fits the definition of Empathetic (feeling with someone) better than Disappointment. Therefore, I must assign a non{-}zero score to Empathetic to capture this nuance of their relationship, acknowledging that his mockery is rooted in a deep, albeit twisted, understanding of her fears.} \\
\\ 
\texttt{Trial \#2: New plan: The previous attempt failed because the scores were identical to the failed attempt, indicating a lack of strategic adjustment despite the analysis. The core issue is that the "Empathetic" score was set to 6, which might be too high given the context of a bitter rivalry where "mockery" is the primary mode of interaction. Remy's final line is a sarcastic jab, not a genuine expression of empathy. Therefore, the "Empathetic" score should be reduced to 0 or 1 to accurately reflect that he feels *mocking* and *disappointed* rather than empathetic. The "Mocking" score should remain high (9) as it is his consistent stance. The "Astonished" score can stay low (3) as he is not surprised by her admission, just amused by it. The "Disappointed" score should be increased (e.g., to 4) because his final line ("And you were always afraid to lose") implies a belief that she is weak or incapable, which is a form of disappointment in her character. The new plan is to recalibrate the scores to reflect a dominant mocking tone with a slight increase in disappointment and a significant decrease in empathy.} \\
\\ 
\texttt{Corrected Scores:} \\
\texttt{Empathetic: 0} \\
\texttt{Mocking: 9} \\
\texttt{Astonished: 3} \\
\texttt{Disappointed: 4} \\
\\ 
\texttt{Trial \#3: New plan: The previous attempts failed because the output format did not strictly adhere to the required structure (specifically, the lack of a "New plan:" header in the final submission block and potential ambiguity in the explanation section). The core task requires predicting Remy's emotion based on the dialogue's climax. Remy's final line ("And you were always afraid to lose, Eilis") is a direct, sarcastic rebuttal to Eilis's admission of fear. This confirms his primary emotion is **Mocking** (9) as he mocks her perceived weakness. However, the prompt asks for the emotion Remy feels *at the end*. While he is mocking, the act of calling her out on her fear implies a sense of **Disappointment** (4) in her inability to rise to the challenge or his own expectations, rather than empathy. The "Astonished" score (3) remains valid as he is surprised by her honesty. The "Empathetic" score is 0 because his tone is hostile. The critical fix is to ensure the final output is a clean, formatted list of scores without extra conversational text, and to explicitly state the reasoning for the "Disappointed" score in the thought process to justify the non{-}zero value, ensuring the strategy accounts for the nuance of his sarcasm masking disappointment.} \\
\\ 
\texttt{Corrected Scores:} \\
\texttt{Empathetic: 0} \\
\texttt{Mocking: 9} \\
\texttt{Astonished: 3} \\
\texttt{Disappointed: 4} \\
\\ 
\\ 
\\ 
\texttt{Your task is to predict the likely emotional responses of a character in this dialogue:} \\
\\ 
\texttt{Eilis: Remy, you're late.} \\
\texttt{Remy: And you're as punctual as ever, Eilis. Still chasing shadows, are we?} \\
\texttt{Eilis: Better shadows than illusions of grandeur, Remy.} \\
\texttt{Remy: Ah, the sting of your wit. I've missed that. } \\
\texttt{Eilis: And I've missed your incessant need to prove something.} \\
\texttt{Remy: Is that so? Well, I suppose we're both here for the same reason then.} \\
\texttt{Eilis: Don't flatter yourself, Remy. Our reasons couldn't be more different.} \\
\texttt{Remy: Oh, do enlighten me, Eilis. } \\
\texttt{Eilis: You're here to feed your ego. I'm here to reclaim what I lost.} \\
\texttt{Remy: And what if I told you I'm here for you, Eilis?} \\
\texttt{Eilis: I'd say you're a bigger fool than I thought.} \\
\texttt{Remy: Well, then let's be fools together, shall we? } \\
\texttt{Eilis: I don't need your charity, Remy.} \\
\texttt{Remy: It's not charity, Eilis. It's competition. } \\
\texttt{Eilis: Is that a challenge, Remy?} \\
\texttt{Remy: If you're not afraid, Eilis.} \\
\texttt{Eilis: Afraid? Of you? Don't make me laugh.} \\
\texttt{Remy: Then let's make a wager, Eilis. } \\
\texttt{Eilis: A wager? And what would that be?} \\
\texttt{Remy: The first to create the Philosopher's Stone.} \\
\texttt{Eilis: And what does the loser get, Remy?} \\
\texttt{Remy: The satisfaction of knowing they gave it their all, Eilis.} \\
\texttt{Eilis: You always were a romantic fool, Remy.} \\
\texttt{Remy: And you were always afraid to lose, Eilis.} \\
\texttt{Eilis: I'm not afraid to lose, Remy. I'm afraid to fail.} \\
\texttt{{[}End dialogue{]}} \\
\\ 
\texttt{At the end of this dialogue, Remy would feel...} \\
\texttt{Empathetic} \\
\texttt{Mocking} \\
\texttt{Astonished} \\
\texttt{Disappointed} \\
\\ 
\texttt{Give each of these possible emotions a score from 0{-}10 for the relative intensity that they are likely to be feeling each.} \\
\\ 
\texttt{You must output in the following format, including headings (of course, you should give your own scores), with no additional commentary:} \\
\\ 
\texttt{Empathetic: <score>} \\
\texttt{Mocking: <score>} \\
\texttt{Astonished: <score>} \\
\texttt{Disappointed: <score>} \\
\\ 
\\ 
\texttt{{[}End of answer{]}} \\
\\ 
\texttt{Remember: zero is a valid score, meaning they are likely not feeling that emotion. You must score at least one emotion > 0.} \\
\\ 
\texttt{Your answer:} \\
            }
        }
    }
     \\ \\

    \theutterance \stepcounter{utterance}  
    & \multicolumn{4}{p{0.6\linewidth}}{
        \cellcolor[rgb]{0.561,0.737,0.561}{
            \makecell[{{p{\linewidth}}}]{
                \texttt{\tiny{[P1$\rangle$GM]}}
                \texttt{Empathetic: 0} \\
\texttt{Mocking: 9} \\
\texttt{Astonished: 3} \\
\texttt{Disappointed: 4} \\
            }
        }
    }
    & & \\ \\

    \theutterance \stepcounter{utterance}  
    & & & \multicolumn{2}{p{0.3\linewidth}}{
        \cellcolor[rgb]{0.9,0.9,0.9}{
            \makecell[{{p{\linewidth}}}]{
                \texttt{\tiny{[GM$|$GM]}}
                \texttt{Empathetic: 0} \\
\texttt{Mocking: 9} \\
\texttt{Astonished: 3} \\
\texttt{Disappointed: 4} \\
            }
        }
    }
    & & \\ \\

    \theutterance \stepcounter{utterance}  
    & & & \multicolumn{2}{p{0.3\linewidth}}{
        \cellcolor[rgb]{0.9,0.9,0.9}{
            \makecell[{{p{\linewidth}}}]{
                \texttt{\tiny{[GM$|$GM]}}
                \texttt{Empathetic: 6} \\
\texttt{Mocking: 7} \\
\texttt{Astonished: 0} \\
\texttt{Disappointed: 2} \\
            }
        }
    }
    & & \\ \\

    \theutterance \stepcounter{utterance}  
    & & & \multicolumn{2}{p{0.3\linewidth}}{
        \cellcolor[rgb]{0.9,0.9,0.9}{
            \makecell[{{p{\linewidth}}}]{
                \texttt{\tiny{[GM$|$GM]}}
                \texttt{game\_result = LOSE} \\
            }
        }
    }
    & & \\ \\

    \theutterance \stepcounter{utterance}  
    & & \multicolumn{4}{p{0.6\linewidth}}{
        \cellcolor[rgb]{0.95,0.95,0.95}{
            \makecell[{{p{\linewidth}}}]{
                \texttt{\tiny{[P1$\leftrightarrow$P]}}
                \texttt{Post{-}episode Reflexion (not sent to GM):} \\
\texttt{New plan: The previous attempts failed because the output format included conversational text ("New plan:", "Corrected Scores:", explanations) instead of strictly adhering to the required format which demands only the four emotion scores with no additional commentary. The strategy must be to output *only* the four lines in the exact format: "Empathetic: <score>", "Mocking: <score>", "Astonished: <score>", "Disappointed: <score>". Based on the dialogue, Remy's primary emotion is Mocking (9) as he consistently belittles Eilis. His final line is a sarcastic jab, so Empathetic is 0. He is slightly surprised by her honesty, so Astonished is 3. He is disappointed in her perceived weakness, so Disappointed is 4. The final output must contain only these four lines.} \\
\\ 
\texttt{Empathetic: 0} \\
\texttt{Mocking: 9} \\
\texttt{Astonished: 3} \\
\texttt{Disappointed: 4} \\
            }
        }
    }
    & \\ \\

\end{supertabular}
}

\end{document}
